\section{Mathematical targets at the paper level}
\label{sec:math-target}

\subsection{Minimal HAETAE-2 background}

The paper describes HAETAE as a module-lattice signature scheme over
\[
  R=\mathbb{Z}[X]/(X^{256}+1),\qquad
  R_q=\mathbb{Z}_q[X]/(X^{256}+1),\quad q=64513
\]
\cite{haetae-spec}. In mode 2, \((k,\ell)=(2,4)\), the small-secret
coefficient bound is \(\eta=1\), and the binary challenge has weight \(\tau=58\).
The verification key is compressed as \(vk=(\mathit{seedA},b_1)\); its mode-2
length in the Jasmin API is 992 bytes.

The core data flow of Figure~8 in the paper, restricted to what this guide
needs, is as follows.

\begin{align*}
  b &= a_{\mathrm{gen}}+A_{\mathrm{gen}}s_{\mathrm{gen}}+e_{\mathrm{gen}},
  & (b_0,b_1)&=(\operatorname{LowBits}_{vk}(b),
                    \operatorname{HighBits}_{vk}(b)),\\
  vk&=(\mathit{seedA},b_1),
  & \mu&=H_{\mathrm{gen}}(\mathit{seedA},b_1,M),\\
  \rho&=H(\operatorname{HighBits}_h(w),
          \operatorname{LSB}(\lfloor y_0\rceil),\mu),
  & c&=\operatorname{SampleBinaryChallenge}_{\tau}(\rho).
\end{align*}

Sign constructs \(\mu\) from the verification-key component stored inside the
secret key, while Verify reconstructs it from the public verification key.
Both then place highbits, LSB, and \(\mu\) in the challenge transcript.
The first implementation-correctness composition point therefore consists of
three concrete equalities, rather than merely saying that ``SHAKE is correct.''

\begin{enumerate}
  \item The packed secret-key prefix read by Sign equals the packed
        verification-key bytes read by Verify.
  \item After absorbing the same pre and message, the first 32 bytes of Sign's
        64-byte \(\mu\) equal Verify's 32-byte \(\mu\).
  \item Absorbing those 32 bytes at position 64 of the same challenge state
        produces equal final states and positions.
\end{enumerate}

The work connected these three equalities to actual generated procedures on
raw/internal paths, tracing the control/memory facts by which actual raw-API
callers reach the inputs, the Sign output frame, and hash equivalence requiring
only the regions read. Directly connecting the two \(\mu\) values stored in
completed traces was frozen at a product/replay semantic boundary.
The main execution path in Weeks~8--15 was the actual signature codec's HBZ
rANS suffix. Coefficient/symbol leaves, tables, pure single steps, the common
byte stack, actual suffix copying, and 3-call harness control were closed.
Week~10 added actual encoder inner-byte writes and successful size bounds.
Week~11 composed exact 0/1/2-byte progress preserving the pure tail in the
inner invariant, generated table word updates, the outer one-symbol transition,
and the final four stores, identifying the complete actual successful encoder
suffix with the pure trace. Week~12 closed exact consumption of that trace by
the actual decoder, and Week~13 closed the actual encoder--suffix-copy--decoder
harness composition. Week~14 composed this core inverse with actual
prepare/apply and production full-HBZ wrapper boundaries. Week~15 proved
nonzero-size success for terminating actual encoder executions on the all-6
symbol sequence produced by canonical all-zero HBZ. Week~16 compiled the KeyGen
matrix/finalizer snapshot, the final sound NTT representation rewrite of
\identifier{output_row}, and the actual harness's two row consequences as
target 77. The required convolution layer was absent from the frozen surface,
so the boundary was fixed at \identifier{STOP-KG-NTT}; limited Sign/Verify
results then completed the 82-target \identifier{PAPER-FROZEN} boundary.
Post-freeze work closed pure NTT/Rq representations and the
actual-(a)/paper-(qj) auxiliary layer, but the synthetic security object's
semantic mismatch makes \identifier{STOP-KG-AVEC-QJ} the latest gate.

\subsection{Functional-correctness target for the mode-2 signature codec}

An actual mode-2 signature has 1474 bytes. The extracted packer/unpacker
layout is as follows.
\[
\underbrace{[0,32)}_{256\text{ challenge bits}}
\;\Vert\;
\underbrace{[32,1056)}_{1024\text{ signed low bytes}}
\;\Vert\;
\underbrace{[1056,1058)}_{\text{two size deltas}}
\;\Vert\;
\underbrace{[1058,1474)}_{\text{hbz/h payload and canonical zero padding}}.
\]
The final region has a total budget of \(416\) bytes, giving
\(1056+2+416=1474\). For a canonical challenge \(c\) and a low-part array
\(z_{\mathrm{low}}\) in the signed-byte range, the closed prefix target is
\[
  \operatorname{UnpackPrefix}
    (\operatorname{PackPrefix}(c,z_{\mathrm{low}}))
  =(c,z_{\mathrm{low}})
\]
This equality is proved as partial correctness for actual generated
procedures. It excludes the two rANS payloads, size metadata, zero padding,
and success and uniqueness of the complete parse. The Week~8--12 results
further divide the internal layers of the first, HBZ payload.

\subsection{Mathematical structure of HBZ leaves, rANS single steps, and byte traces}

A mode-2 HBZ input consists of \(1024\) coefficients \(h_i\in[-6,7)\).
The actual prepare procedure maps them to
\[
  s_i=h_i+6\in\{0,\ldots,12\}
\]
The actual apply procedure reconstructs \(s_i-6\) as a 32-bit word. Thus
terminating composition of the two actual leaves on a canonical HBZ array satisfies
\[
  \operatorname{Apply}(\operatorname{Prepare}(h))[0..1024)=h[0..1024)
\]
The prepare theorem separately preserves the symbol tail of capacity
\(2048\), and the composition postcondition preserves the decoded-coefficient
tail. This is leaf-level partial correctness, with no claim about rANS buffer
success or size.

The cumulative starts and frequencies of the 13 symbols are respectively
\[
\begin{aligned}
(c_s)_{s=0}^{12}
  &=(0,1,2,3,8,66,312,710,957,1016,1021,1022,1023),\\
(f_s)_{s=0}^{12}
  &=(1,1,1,5,58,246,398,247,59,5,1,1,1),
\end{aligned}
\]
Here \(f_s>0\), \(\sum_s f_s=1024\), and the intervals
\([c_s,c_s+f_s)\) partition \([0,1024)\). Define the pure rANS steps by
\[
\begin{aligned}
  E_s(x)&=\left\lfloor\frac{x}{f_s}\right\rfloor 1024
           +c_s+(x\bmod f_s),\\
  D_s(y)&=\left\lfloor\frac{y}{1024}\right\rfloor f_s
           +(y\bmod 1024)-c_s
\end{aligned}
\]
The current proof establishes \(D_s(E_s(x))=x\) for \(x\ge0\). It also proves
that the fast step computed by the actual mode-2 reciprocal/bias expression
equals \(E_s(x)\) when \(0\le s<13\) and
\(1\le x<\mathsf{hbz\_xmax}(s)\). The literal
\identifier{esyms}/\identifier{dsyms_words}/\identifier{symbol_words} arrays
in the actual packing/unpacking extractions satisfy this certificate.

Week~9 fixed the pure byte trace linking single steps as follows. Let the
symbols be \(s_0,\ldots,s_{n-1}\), with \(n=1024\), and set
\(x_n=L=2^{23}\). In encoder order \(j=n-1,\ldots,0\),
\[
\begin{aligned}
  \ell_j&=\operatorname{renorm\_len}
      (x_{j+1},\mathsf{hbz\_xmax}(s_j))\in\{0,1,2\},\\
  r_j&=\left\lfloor x_{j+1}/256^{\ell_j}\right\rfloor,\\
  x_j&=E_{s_j}(r_j).
\end{aligned}
\]
Let \(\beta_j\) be the normalization-byte segment in the order read by the
actual decoder. The compiled readback theorem gives
\[
  \operatorname{ReadBytes}(r_j,\beta_j)=x_{j+1}
\]
In particular, when emitting two bytes, the actual encoder writes the low
byte first while decrementing its cursor, so the final forward segment is
\([\mathsf{high},\mathsf{low}]\). The common trace is therefore
\[
  c_0=4,\qquad c_{j+1}=c_j+|\beta_j|,
  \qquad
  B=\operatorname{LE32}(x_0)\Vert\beta_0\Vert\cdots\Vert\beta_{n-1}
\]
Here \(\operatorname{ParseLE32}\) restores the first state; each
\(D_{s_j}\) obtains \(r_j\) from \(x_j\), then reads \(\beta_j\) to recover
\(x_{j+1}\). The pure induction theorem preserves
\[
  2^{23}\le x_j<2^{31}
\]
at every boundary state.

These pure results alone do not imply refinement of the actual encoder and
decoder outer loops. Compiled results also include the Hoare theorem that
actual \identifier{__copy_encoded_suffix} copies the encoder suffix pointwise
to a decoder prefix while preserving the remaining tail, and the unary
harness-control theorem that performs or skips copying and decoding according
to the actual encoder's returned \(bad\).

\subsection{Week~10 mathematical boundary of the actual encoder}

For a symbol array \(a\) of capacity 2048, define the list used by mode 2 and
its suffix by
\[
  S(a)=[\operatorname{uint}(a[0]),\ldots,
        \operatorname{uint}(a[1023])],\qquad
  S_i(a)=\operatorname{drop}_i(S(a))
\]
These are exactly \identifier{symbol_list_of_array} and
\identifier{symbol_suffix} in the code. Read the array segment and frame conditions as
\[
\begin{aligned}
  \operatorname{Segment}(A,p,b)
    &\Longleftrightarrow
      0\le p\;\land\;p+|b|\le2048\;\land\\
    &\hspace{2.8em}\forall k\in[0,|b|),\;
      \operatorname{uint}(A[p+k])=b[k],\\
  \operatorname{PrefixFrame}(A_0,A_1,p)
    &\Longleftrightarrow
      \forall k\in[0,p),\;A_0[k]=A_1[k]
\end{aligned}
\]
The compiled array/list bridges give \(S_{1024}(a)=[]\) and
\(S_i(a)=\operatorname{uint}(a[i])::S_{i+1}(a)\). They also isolate the
one-symbol trace recurrence as follows.
\[
\begin{aligned}
  T(s::t).\mathsf{state}
    &=E_s(\operatorname{Norm}(T(t).\mathsf{state},s)),\\
  T(s::t).\mathsf{bytes}
    &=\operatorname{NormBytes}(T(t).\mathsf{state},s)
      \Vert T(t).\mathsf{bytes}.
\end{aligned}
\]

The word step modeling the actual table's four words, 64-bit reciprocal
multiplication and high-half extraction, shift ladder, and bias satisfies
\[
\begin{gathered}
  0\le s<13,\qquad
  1\le\operatorname{uint}(x)<\mathsf{hbz\_xmax}(s),\\
  \operatorname{ActualWordStep}(x,s)
   =\operatorname{W32}
      (\operatorname{FastStep}(\operatorname{uint}(x),s))
\end{gathered}
\]
Combined with the preceding fast-step theorem, the right side is the W32
representation of mathematical \(E_s\). This is a word-level equality, not a
Hoare theorem for the generated outer loop.

The generated \identifier{_rans_encode} inner normalization loop carries the
following local invariant, with witnesses for entry state \((x_0,off_0)\),
threshold \(x_{\max}\), and step \(k\).
\[
\begin{gathered}
  4\le off_0\le1024,\qquad 0\le k\le4,\qquad
  1\le\operatorname{uint}(x_{\max}),\\
  off=off_0-k,\qquad
  x=\operatorname{PhaseState}(x_0,k),\\
  \operatorname{Segment}
    (encp,off_0-k,\operatorname{PhaseBytes}(x_0,k)),\qquad
  \operatorname{PrefixFrame}(before,encp,off_0-k).
\end{gathered}
\]
It directly proves that each actual \identifier{set8} step prepends the low
byte to the list, the W64 cursor does not underflow, and the unwritten prefix
is preserved. However, \(before\) is the array at entry to that inner loop,
so the invariant does not yet describe the complete accumulated outer trace.
The bound \(0\le k\le4\) is conservative, derived from W32 finiteness;
identification with the pure trace's exact 0/1/2-byte length remains open.

Consequently, two partial-correctness theorems for the actual generated
procedure establish
\[
\begin{aligned}
 bad'\ne0\quad\lor\quad
   &(bad'=0\land 0\le off'\le1020),\\
 bad'=0\quad\Longrightarrow\quad
   &4\le1024-off'\le1024
\end{aligned}
\]
The second formula bounds encoded size on success. The conclusions do not
include reachability or termination of a successful path, or success on input whose symbols are all 6 (all-6 input).

\subsection{Week~11 closure of the actual encoder trace}

Week~11's tail-aware inner invariant carries the already written pure tail
\(t\), in addition to the local bytes.
\[
\begin{aligned}
 0\le k\le\ell(x_0,s),\qquad
 off&=off_0-k,\\
 \operatorname{Segment}\bigl(encp,off_0-k,
   \operatorname{Written}(x_0,s,k)\Vert T(t).\mathsf{bytes}\bigr),\qquad
 &\operatorname{PrefixFrame}(enc_0,encp,off_0-k).
\end{aligned}
\]
Here \(\ell(x_0,s)\in\{0,1,2\}\) is the pure normalization length. Proving
\(k=\ell(x_0,s)\) when the actual generated guard becomes false closes the gap
between Week~10's conservative \(0..4\) phase and the pure trace.

The successful-path outer invariant is
\[
\begin{gathered}
  0\le i\le1024,\qquad
  x=T(S_i(a_0)).\mathsf{state},\qquad 4\le off,\\
  off+|T(S_i(a_0)).\mathsf{bytes}|=1024,\\
  \operatorname{Segment}(encp,off,T(S_i(a_0)).\mathsf{bytes}),\qquad
  \operatorname{PrefixFrame}(enc_0,encp,off)
\end{gathered}
\]
Dedicated word theorems fold generated table reads and protection operations,
reciprocal multiplication, and the shift ladder to recover the one-symbol
transition. At \(i=0\), the four actual little-endian stores are composed.
A terminating actual \identifier{_rans_encode} thus either returns failure
or satisfies the following postcondition on success.
\[
\begin{aligned}
 bad'\ne0\quad\lor\quad\bigl(bad'=0
 &\land 0\le off'\le1020
 \land 4\le1024-off'\le1024\\
 &\land\operatorname{Segment}(encp',off',
       \operatorname{trace\_bytes}(S(a_0)))\\
 &\land off'+|\operatorname{trace\_bytes}(S(a_0))|=1024\\
 &\land\operatorname{PrefixFrame}(enc_0,encp',off')\bigr).
\end{aligned}
\]
\identifier{actual_rans_encode_trace_closure} and
\identifier{actual_rans_encode_trace_refinement} prove this statement for the
actual generated procedure. Thus \claim{OBL-RANS-ENCODE-REFINEMENT} is
\statusproved as success-conditioned partial correctness. Success is not a
precondition, but the theorem does not prove termination or reachability of \(bad'=0\).

\subsection{Week~12 closure of the actual decoder trace}

For the expected symbol array \(a\) of capacity 2048, define
\[
  s:=S(a),\qquad B:=\operatorname{trace\_bytes}(s),\qquad n:=|B|
\]
The actual decoder's exact-trace input requires
\identifier{mode2_hbz_symbol_stream}(a), \(4\le n\le1024\), pointwise
agreement of the input-buffer segment \([0,n)\) with \(B\), state fields
\[
  (\mathsf{count},\mathsf{size},m)=(1024,n,13)
\]
and bindings for the actual \identifier{symbol_words}/\identifier{dsyms_words}
tables. It also explicitly requires pointwise reads within each normalization
segment to equal the corresponding pure list elements. This input predicate
contains no decoder success, output-symbol equality, or post-state assumption.

Defining the reference state, cursor, and segment by
\[
\begin{aligned}
  x_i&:=\operatorname{decoder\_state\_at}(s,i),&
  c_i&:=\operatorname{decoder\_cursor}(s,i),\\
  \beta_i&:=\operatorname{decoder\_segment\_at}(s,i)
\end{aligned}
\]
the compiled boundary theorem gives
\[
  c_0=4,\qquad c_{i+1}=c_i+|\beta_i|,\qquad
  c_{1024}=n,\qquad x_{1024}=2^{23}
\]
The actual inner loop carries \(k=\mathsf{off}-c_i\) and the partially replayed
state, consuming precisely the next byte or stopping according to segment
length \(|\beta_i|\in\{0,1,2\}\). Separate theorems connect concrete packed-table
lookup and W32/W64 operations to selection of the same symbol and mathematical
decoder step.

Consequently, \identifier{actual_rans_decode_trace_refinement} proves the
following for every terminating execution of the actual generated
\identifier{RansDecodeTarget.M._rans_decode}:
\[
\begin{aligned}
  bad_{\mathrm{dec}}&=0,&
  off_{\mathrm{dec}}&=n,\\
  decoded'[0..1024)&=a[0..1024),&
  decoded'[1024..2048)&=decoded_0[1024..2048)
\end{aligned}
\]
The internal fact \(x_{1024}=2^{23}\) excludes the actual final rejection
check, but is not exposed in the generated ABI result. Thus
\claim{OBL-RANS-DECODE-REFINEMENT} is \statusproved as exact-trace partial
correctness. Termination, losslessness, and actual encoder-success reachability
are not conclusions of this theorem.

\subsection{Week~13 closure of actual rANS core composition}

Week~13 sequentially composed the preceding three individual theorems inside
actual \identifier{Mode2RansActualHarness.run}. This harness directly calls the
fixed extractions' \identifier{_rans_encode},
\identifier{__copy_encoded_suffix}, and \identifier{_rans_decode}, in that order.
\identifier{Mode2RansCoreCompositionBridge.ec} constructs the copied buffer's
exact trace from the encoder's \identifier{segment_matches} and actual copying's
\identifier{slice_eq}, then transports it to every pointwise decoder read and
the constructed state fields \((1024,n,13)\). It also derives from the actual
successful size range that W64 subtraction equals the integer trace size
without wraparound.

The only preconditions of \identifier{actual_rans_encode_copy_decode_inverse}
are initial argument bindings and \identifier{mode2_hbz_symbol_stream}. In
particular, encoder success and copy/decoder results are not preconditions.
For every terminating execution, the conclusion is approximately
\[
\begin{aligned}
  &\bigl(bad_{\mathrm{enc}}\ne0
    \land\neg\mathsf{decoder\_ran}
    \land decoded'=decoded_0\bigr)\\
  {}\lor{}&\bigl(bad_{\mathrm{enc}}=0
    \land\mathsf{decoder\_ran}
    \land bad_{\mathrm{dec}}=0
    \land off_{\mathrm{dec}}=\mathsf{size}\\
    &\land decoded[0..1024)=symbols[0..1024)
    \land\mathsf{frames}\bigr)
\end{aligned}
\]
On failure the decoder state also retains its initial value. On success,
the decoder input fields, actual encoder trace suffix, and initial-prefix frame
are preserved as well. Thus \claim{OBL-RANS-CORE-INVERSE} is \statusproved as
success-conditioned Hoare partial correctness.

This is a theorem about a proof-authored harness directly calling three
actual generated procedures, not yet a wrapper theorem for
production \identifier{_encode_hb_z1_full}/\identifier{_decode_hb_z1_full}.
Canonical \(h\) alone also does not imply buffer success, so the full HBZ
wrapper target must still expose at least the actual failure branch as follows:
\[
  \mathsf{size}=0\quad\lor\quad
  \bigl(\mathsf{size}\ne0\land
        \mathsf{bad}=0\land
        \mathsf{decoded}[0..1024)=h[0..1024)\bigr)
\]
Transport of the full HBZ parent's frame and size, encoder/decoder
termination, and an actual success witness for all-6 input had not yet been compiled.

\subsection{Week~14 closure of production full-HBZ wrappers}

Week~14 connected the focused full-HBZ path calling
\identifier{_encode_hb_z1_full} followed by \identifier{_decode_hb_z1_full}
to the production SignaturePack/Unpack path, showing that they expose the same
actual procedure semantics. \identifier{Mode2HbzInternalBoundaries.ec}
establishes exact equivalence between focused/full extractions of prepare,
rANS encoder, rANS decoder, and apply. It provides adapters from canonical
coefficients to a 13-symbol stream and from recovered symbols back to apply premises.

The only semantic premises of
\identifier{actual_hbz_full_encode_decode_inverse_mode2} are initial argument
bindings and \identifier{canonical_hbz_mode2}. Transported to the production
boundary, \identifier{signature_pack_unpack_hbz_full_inverse_mode2} proves one
of the following two branches for every terminating execution.
\[
\begin{aligned}
  \mathsf{size}=0
    &\Longrightarrow \neg\mathsf{decoder\_ran}
      \land encoded'=encoded_0
      \land decoded'=decoded_0
      \land bad'=bad_0,\\
  \mathsf{size}\ne0
    &\Longrightarrow 4\le\mathsf{size}\le1024
      \land\mathsf{decoder\_ran}
      \land bad'=0\\
    &\hspace{4.2em}\land\;
      decoded'[0..1024)=hbz_0[0..1024)
      \land\mathsf{coefficient/encoded\ frames}.
\end{aligned}
\]
The success branch also includes the \identifier{prepared_symbols} witness
produced by actual prepare and equality of its \identifier{mode2_trace_bytes}
with the returned encoded segment.
\identifier{signature_pack_unpack_hbz_full_actual_exact} first proves exact
equality of the production and focused harnesses' result tuples and
\identifier{decoder_ran}, so generated-procedure semantics are not assumed as a new axiom.

The following obligation is therefore closed as success-conditioned Hoare partial correctness.
\begin{center}
  \claim{OBL-SIG-HBZ-ENCODE-DECODE}: \statusproved
\end{center}
However, this statement does not eliminate the \(\mathsf{size}=0\) branch.
An actual success witness for all-zero coefficients (all-6 symbols after
prepare) and encoder/decoder termination remain separate obligations. Before
those are closed, unconditional round trips or \(\delta_{\mathrm{Encoding}}=0\)
cannot be concluded.

\subsection{Week~15 actual success postcondition for fixed all-6 input}

Week~15 closed the previously isolated fixed-input edge with two authored
theories. \sourcepath{Mode2RansAllSix}\allowbreak\sourcepath{Budget.ec}
unfolds every 32-bit coefficient read of bytewise-zero-initialized
\identifier{zero_hbz} to 0, establishing canonicality, and proves that the
first 1024 symbols produced by actual prepare are all 6.

The frequency and threshold for symbol 6 are
\[
  \mathsf{freq}(6)=398,\qquad
  \mathsf{xmax}(6)=2^{21}\cdot398=834666496
\]
When the state is below \(2^{31}\), one symbol emits at most one
normalization byte, and the first four pure states
\[
8388608\to21582496\to55528910\to142868116\to367580518
\]
emit no bytes. Thus the normalization portion of 1024 all-6 symbols uses at
most 1020 bytes. Including 4 bytes for final-state serialization, the complete trace satisfies
\[
  4\le |\mathsf{mode2\_trace\_bytes}|\le1024
\]

The strengthened actual encoder-closure theorem preserves
\[
  bad\ne0\Longrightarrow
  1020<|\mathsf{normalization\ trace}|
\]
This contradicts the pure bound above. Hence
\identifier{actual_rans_encode_all_six_success} concludes \(bad=0\), exact trace
and size, and the prefix frame for terminating actual \identifier{_rans_encode}
executions. Lifting this to full-HBZ and production boundaries,
\identifier{signature_pack_unpack_hbz_zero_success_mode2} proves, on fixed
\identifier{zero_hbz} input,
\[
  size\ne0,\quad 4\le size\le1024,\quad
  \mathsf{decoder\_ran},\quad bad_{\mathrm{dec}}=0,\quad
  decoded[0..1024)=0
\]
and both tail frames.

This makes \claim{OBL-RANS-ACTUAL-SUCCESS-WITNESS} \statusproved as Hoare
partial correctness for fixed all-6 input. A \(\operatorname{hoare}\) statement,
however, does not establish that the execution exists or terminates. Fixed-input
losslessness, \(\operatorname{phoare}=1\), success on general canonical inputs,
and encoder/decoder termination remain separate obligations.

\subsection{Week~16 actual KeyGen matrix/finalizer snapshot}

The actual boundary of the paper's KeyGen equation was restricted to two
fixed helpers, rather than opening the entire retry/API layer at once.
\identifier{ActualM23MatrixFinalizeSnapshot.run} calls actual
\identifier{_kp_m23_matrix} with mode-2 dimensions \((2,3)\), stores its returned
array as \(p=\mathsf{pre\_bp}\), then calls actual
\identifier{_keypair_finalize_m23} on 512 words.

For each active coefficient, let
\[
  r=(p_i+s_{2,i}+a_i)\bmod q,\qquad
  b_{0,i}=\operatorname{Low}(r),\qquad
  s'_{2,i}=s_{2,i}-b_{0,i}
\]
The compiled semantic theorem then gives
\[
  r=2b_{1,i}+b_{0,i},\qquad -1\le b_{0,i}\le1
\]
and the snapshot-scoped congruence
\[
  2(a_i-2b_{1,i})+2p_i+2s'_{2,i}\equiv0\pmod{2q}
\]
These are the actual finalizer's KG-2-style low/high decomposition and the
adjusted \(s_2\) component of KG-4. The matrix/NTT representation and all three
array-tail frames are preserved as well.

Crucially, \(p\) is the \identifier{KeygenM23ArithmeticSpec.output_row}
snapshot returned by the actual matrix helper.
\identifier{output_row_from_mode2_ntt_words} compiles the final representation
edge expanding this value as
\[
 \operatorname{array256\_mont}
 \bigl(\operatorname{full\_invntt}
   (\operatorname{pointwise\_row\_words}(m,v,row))\bigr)
\]
\identifier{actual_m23_matrix_snapshot_rows_explicit} transports it to the
two active-row postconditions of the actual two-call harness. This expression
is not renamed to the paper's \(A_{\mathrm{gen}}s_{\mathrm{gen}}\). At paper
freeze, the first missing leaf identified by the audit was odd-root
orthogonality/full-NTT spectral action with the actual Montgomery factor:
identifying inverse NTT of the pointwise product with \identifier{Rq.\&*}
negacyclic convolution. A separate adapter was then required from
\identifier{Rq.poly} arrays to \identifier{HAETAE_Algebra.poly_mul} integer
lists. Hence, in the frozen claim matrix, faithful KG-1, KG-3, complete KG-4, and
\[
  A s=qj\pmod{2q}
\]
were not proved. The frozen \identifier{STOP-KG-NTT} decision fixed KeyGen
at KG-2 and the finalization boundary. Direct calls to the three actual Sign
helpers and accepted-branch control were subsequently compiled, but the following theorem is absent.
\begin{center}
  \small\identifier{sf_challenge_mode2_highbits_lsb_sampleinball_correct}
\end{center}
The boundary was therefore frozen as \identifier{STOP-SIGN-CHAL-MODE2}.
The S-1--S-7 headline theorem is unproved. The subsequent MINCORE-VERIFY work
preserved helper-local machine-word and control results.

After paper freeze, \identifier{rq_mul_coeff_foldr_to_bigi},
\identifier{full_ntt_montgomery_}\allowbreak\identifier{spectral_action},
general row products, and the Rq--HAETAE multiplication-representation
adapter were closed in a separate auxiliary tree. Actual mode-2
\identifier{avec} was also connected to the paper's \(a\) through its
SHAKE128/rejection semantics, and the paper's \(qj\) term was added
independently to the snapshot equation. However, the existing security model's
synthetic \identifier{haetae_mode23_qj_vector} was proved different from these
objects, yielding the latest decision \identifier{STOP-KG-AVEC-QJ}. KG-1--KG-4
and the paper equation above therefore remain unpromoted. For Verify, the
first procedure blockers of the combined theorem are actual 2-by-4
\identifier{verify_matrix_ntt_acc_mode2_cols4_correct} and
\identifier{verify_crt_freeze_mode2_word_exact}, rather than the pure NTT foundation.

\subsection{Top-level functional-correctness targets}

\sourcepath{easycrypt/specs/TopLevelGoals.ec} models the implementation and
the paper's probabilistic algorithms as abstract operations and defines the following targets.

\begin{align*}
  \mathsf{FC}_{\mathrm{KG}}
    &:\quad \mathsf{JKeyGen}_2=\mathsf{PaperKeyGen}_2,\\
  \mathsf{FC}_{\mathrm{Sign}}
    &:\quad \forall sk,M,\;
       \mathsf{JSign}_2(sk,M)=\mathsf{PaperSign}_2(sk,M),\\
  \mathsf{FC}_{\mathrm{Verify}}
    &:\quad \forall vk,M,\sigma,\;
       \mathsf{JVerify}_2(vk,M,\sigma)
       =\mathsf{PaperVerify}_2(vk,M,\sigma).
\end{align*}

These formulas are specification interfaces to be implemented, not declarations
of completed proofs. The current prefix and transcript theorems are
deterministic subtheorems required by these equalities.

\subsection{Top-level security-composition target}

The target comparing the implementation EUF-CMA advantage with that of the paper scheme is
\begin{equation}
\label{eq:implementation-bound}
  \operatorname{Adv}^{\mathrm{euf\mbox{-}cma}}_{\mathrm{Jasmin}}
  \le
  \operatorname{Adv}^{\mathrm{euf\mbox{-}cma}}_{\mathrm{paper}}
  +\delta_{\mathrm{KeyGen}}
  +q_{\mathrm{Sign}}\delta_{\mathrm{Sign}}
  +\delta_{\mathrm{Verify}}
  +\delta_{\mathrm{Encoding}}
\end{equation}
At the paper level, the further target is
\begin{equation}
\label{eq:paper-bound}
  \operatorname{Adv}^{\mathrm{euf\mbox{-}cma}}_{\mathrm{paper}}
  \le
  \operatorname{Adv}_{\mathrm{MLWE}}
  +\operatorname{Adv}_{\mathrm{BST\mbox{-}MSIS}}
  +\varepsilon_{\mathrm{ROM}}
  +\varepsilon_{\mathrm{Rejection}}
  +\varepsilon_{\mathrm{Fork}}
\end{equation}
Equations \cref{eq:implementation-bound,eq:paper-bound} also remain
\statusspecified or conditional connections. The final cryptographic-assumption
interfaces may be limited to HAETAE-2 MLWE, BST-MSIS, and ROM, but the concrete
game instances and reduction theorems are not all connected yet.
